\documentclass[11pt,a4paper]{article}
\usepackage[T1]{fontenc}
\usepackage[utf8]{inputenc}
\usepackage{lmodern}
\usepackage[a4paper,margin=26mm,headheight=15pt]{geometry}
\usepackage{amsmath,amssymb,amsthm,mathtools}
\usepackage{booktabs,tabularx,longtable,array}
\usepackage{enumitem,microtype,xcolor,xurl}
\usepackage[most]{tcolorbox}
\usepackage{fancyhdr,needspace}
\usepackage[colorlinks=true,linkcolor=blue!45!black,citecolor=blue!45!black,urlcolor=blue!45!black]{hyperref}
\definecolor{ink}{HTML}{16324F}
\definecolor{pale}{HTML}{F0F5F8}
\hypersetup{pdftitle={Replacing primitive labels with norm tests},pdfauthor={ChatGPT},pdfsubject={A proved local simplification proposal for the Lean RH and Dirichlet GRH proof}}
\pagestyle{fancy}
\fancyhf{}
\fancyhead[L]{\small\sffamily Replacing primitive labels with norm tests}
\fancyhead[R]{\small 27 November 2026}
\fancyfoot[C]{\small\thepage}
\setlength{\parindent}{0pt}
\setlength{\parskip}{5.5pt plus 1pt}
\setlist{itemsep=3pt,topsep=5pt,leftmargin=*}
\setcounter{tocdepth}{2}
\newcommand{\F}{\mathbf F}
\newcommand{\ord}{\operatorname{ord}}
\newcommand{\Nm}{\operatorname{Norm}}
\newcommand{\Qual}{\mathcal Q}
\newcommand{\src}[1]{\nolinkurl{#1}}
\newcommand{\commit}{\texttt{89d12dd40424f1217fd818d84877b3da52396a3b}}
\newcommand{\repobase}{https://github.com/VladimirReshetnikov/ProveIt/blob/89d12dd40424f1217fd818d84877b3da52396a3b/Analysis/GeneralizedRiemannHypothesis}
\newcommand{\sourcefile}[2]{\href{\repobase/Lean/grh/src/#1}{#2}}
\newtheorem{theorem}{Theorem}[section]
\newtheorem{lemma}[theorem]{Lemma}
\newtheorem{corollary}[theorem]{Corollary}
\theoremstyle{definition}
\newtheorem{definition}[theorem]{Definition}
\theoremstyle{remark}
\newtheorem{remark}[theorem]{Remark}
\newtcolorbox{keybox}{colback=pale,colframe=ink,boxrule=.5pt,arc=1mm,left=3mm,right=3mm,top=2mm,bottom=2mm,breakable}
\newcolumntype{Y}{>{\raggedright\arraybackslash}X}
\begin{document}
\begin{center}
{\color{ink}\sffamily\bfseries\LARGE Replacing primitive labels\\[3pt] with norm tests}\par
\vspace{8pt}
{\large A local simplification of the Lean proof of RH\\and GRH for Dirichlet $L$-functions}\par
\vspace{7pt}
{27 November 2026}\par
\end{center}

\begin{abstract}
The proof selects a distinguished finite-field polynomial whose root generates
the entire multiplicative group of an extension. The downstream arguments use
only a fraction of that information: irreducibility that survives every
prime-power map, invariance under certain scalar twists, and divisibility of
the order of a common root power. This report replaces the primitive-root
condition by explicit tests on the norm and one exclusion from a subgroup.
The replacement has an exact elementary count. For the parameters already
present in the source, at least $q^r/(4r)$ monic polynomials qualify, and the
bad-type estimate needs $f\ge4rL$ in place of
$f\ge2r(2r\log q+1)L$. Complete proofs, an audit of the consumers and prior
proposals, reproducible exact checks, and a formalization plan are supplied.
The local mathematical result is proved here; its integration into Lean has
not been implemented or compiled.
\end{abstract}

\begin{keybox}
\textbf{Recommendation.} Change the distinguished witness in
\src{goodTuples}, preserving the allowed-polynomial condition and the
\src{IsCommonRootPower} output interface. Replace the primitive-polynomial
totient estimate by a direct count. This is independent of the existing
proposals concerning spanning, torsion interpolation, quadratic twists, and
analytic cutoffs. It is a simplification of one local part of the proof,
with a stronger local bound; it does not replace the theta-lift or Selmer
arguments.
\end{keybox}

\textbf{Source snapshot.} All repository claims and the nonduplication audit
refer to ProveIt commit \commit. The accompanying source manifest pins each
audited file by its Git blob and SHA-256. The scope of ``GRH'' in this report
is the theorem for Dirichlet $L$-functions proved by the linked package.

\textbf{Notation warning.} The integer $T$ below is the source's
\texttt{P.q\textsubscript{2}\ D.towerExp}, the order of a local character. It is \emph{not}
the auxiliary level denoted by $N$ elsewhere in the source. The small examples
in the executable checks call this same integer \src{N}.

\newpage
\begingroup
\small
\setlength{\parskip}{0pt}
\tableofcontents
\endgroup
\newpage

\section{What the proof does, and where the excess hypothesis occurs}
\label{sec:context}

\subsection{The comparison that drives the proof}

The readable source organizes the arithmetic argument around three inputs.
Short integral forms arise from a chain of theta lifts and congruence
estimates. A spanning argument produces sufficiently many independent
translates of those forms. A determinant or height estimate says that
independent integral forms cannot all be too short. The main terms on the
two sides agree, so comparison leaves an upper bound for a weighted prime
remainder. This is the role of Theorems A, B, and C in the guide and of
\src{star_inequality} in the source~\cite{guide,mainineq}.

More explicitly, the implemented terminal arithmetic inequality is
\begin{equation}
 m\,\mathcal R\le f a (\log b)^5,
 \label{eq:main}
\end{equation}
where $\mathcal R$ is \src{D.totalRemainder}. The proof compares an upper
bound on a selected family of forms with a lower bound on its integral
volume. The common main amount cancels, while rank loss and integral-index
cost are absorbed into the right side of~\eqref{eq:main}. Suitable data,
with chain length essentially of square-root size, turn this into an
eventual one-sided bound
\[
 R_{q_0}(x)\le x^{-1/2+\epsilon}
 \qquad(\epsilon>0).
\]
The analytic endgame combines this bound with an appropriate quadratic
twist and a Landau-type pole argument to rule out zeros to the right of
$1/2$~\cite{onesided,guide}. The existing proposals already change several
of these outer steps; this report leaves them available for incorporation.

The local projector sits earlier in the argument. At auxiliary places,
monic squarefree polynomials label finite general-linear-group
representations. A tuple is good when every entry is allowed and at least
one entry is a distinguished irreducible polynomial. The source presently
requires that distinguished polynomial to be primitive: a root must have
order $q^r-1$~\cite{labels}. These good tuples define the projector used
both in the spanning estimate and in the congruence construction.

\subsection{What the distinguished label is actually used for}

There are three substantive requirements after the source chooses a
primitive label.
\begin{enumerate}
\item \textbf{An irreducible local representation.} In the multiplicity
argument the relevant consumer keeps the label and irreducibility
conjuncts and discards the full-order conjunct~\cite{multiplicities}.
\item \textbf{Stability under normalization.} Multiplying all roots by a
base-field $T$-th root of unity must preserve the distinguished condition.
The current proof establishes this for primitive roots by a general
order-of-a-product lemma~\cite{labels,rows}.
\item \textbf{A common root power with sufficient order.} For the scaled
polynomial, all roots must have a common power $y$. If $y^x=1$, the source
needs $T\mid jx$. In addition, the Frobenius conjugates of every root must
remain distinct after raising to $\ell^u$, for every prime $\ell$ and every
$u\ge0$~\cite{rows}.
\end{enumerate}

Primitive roots supply all three properties. They also supply much more:
every prime-power factor of $q^r-1$ occurs at its full exponent. The proof
does not use that additional information after the interface above.
Counting only primitive roots introduces Euler's totient and then a
logarithmic lower bound for the totient. The proposed replacement proves
the required properties directly and counts the resulting elements
exactly.

\subsection{Status of the starting proof and of this analysis}

The linked readable source has 803 Lean files under \src{src/}. The guide's
source checker passes against the inspected checkout: 253 declaration
anchors and 83 Git blobs match. The original proof's independent
verification is documented by Andrew Sutherland: the shorter editions
compile, their five public theorems use only \src{propext},
\src{Classical.choice}, and \src{Quot.sound}, and the reported kernel and
external-checker runs accept the results~\cite{sutherland}.

This work inspected the source interfaces and mathematical argument; it
did not repeat that complete formal verification. No Lean compiler or
kernel run was performed for the proposed change. The distinction matters:
a proved replacement lemma plus an interface audit is a concrete
refactoring proposal, but it is not a newly certified version of the
repository's terminal theorem.

\section{The replacement theorem}
\label{sec:theorem}

Let $q$ be a prime power. Throughout Sections~\ref{sec:theorem}--\ref{sec:count}
assume
\begin{equation}
 r=3^a,\qquad a\ge1,\qquad T>0,\qquad 3\mid T,
 \qquad T^2\mid q-1.
 \label{eq:hyp}
\end{equation}
For each prime $p\mid T$, write $n_p=v_p(T)$. Set
\begin{equation}
 m=q-1,\quad M=q^r-1,\quad E=\frac{M}{m},\quad
 B=q^{r/3},\quad H=3(B-1),\quad C=\frac{B^2+B+1}{3}.
 \label{eq:notation}
\end{equation}
In particular $M=HC$. The hypotheses imply $q\equiv1\pmod9$ and therefore
$B\equiv1\pmod3$; every displayed exponent is an integer.

\begin{definition}[Qualifying roots and polynomials]
For $\alpha\in\F_{q^r}^{\times}$, put $y=\Nm_{\F_{q^r}/\F_q}(\alpha)=\alpha^E$.
Say that $\alpha$ qualifies when
\begin{align}
 y^{m/3}&\ne1, &&\label{eq:test3}\\
 y^{m/p^{n_p}}&\ne1 &&\text{for every prime $p\mid T$, $p\ne3$,}
 \label{eq:testp}\\
 \alpha^H&\ne1. &&\label{eq:outside}
\end{align}
A monic polynomial qualifies when it is the minimal polynomial over $\F_q$
of a qualifying root. Equivalently, one may require that it be a label,
irreducible of degree $r$, with one root satisfying these tests; the theorem
below shows that this formulation defines the same class.
\end{definition}

The first two lines are tests in the base field. The last line excludes a
single subgroup of $\F_{q^r}^{\times}$, of order $H$. It is not merely the
maximal proper subfield's multiplicative group: that group has order
$B-1$, whereas the excluded subgroup has order $3(B-1)$.

\begin{theorem}[Sufficient properties and exact count]
\label{thm:replacement}
Under~\eqref{eq:hyp}, the qualifying roots have the following properties.
\begin{enumerate}[label=\textup{(\roman*)}]
\item For every prime $\ell$ and every integer $u\ge0$,
\[
 [\F_q(\alpha^{\ell^u}):\F_q]=r.
\]
In particular $\alpha$ has degree $r$, and the $r$ conjugates
$(\alpha^{q^i})^{\ell^u}$, $0\le i<r$, are pairwise distinct.
\item The tests are invariant under Frobenius and under
$\alpha\mapsto\lambda\alpha$, where $\lambda\in\F_q^{\times}$ and
$\lambda^T=1$.
\item $T\mid\ord(y)$. All roots of the minimal polynomial of $\alpha$
have the same $E$-th power $y$. Consequently $y^x=1$ implies $T\mid x$,
and hence $T\mid jx$ for every integer $j$.
\item Every primitive element of $\F_{q^r}^{\times}$ qualifies.
\item Define
\begin{equation}
 \delta_T=\frac23\prod_{\substack{p\mid T\\p\ne3}}
                  (1-p^{-n_p}).
 \label{eq:delta}
\end{equation}
The exact number of qualifying roots is
\begin{equation}
 A(q,r,T)=\delta_T\bigl(q^r-3q^{r/3}+2\bigr),
 \label{eq:roots}
\end{equation}
and the exact number of qualifying monic polynomials is
\begin{equation}
 P(q,r,T)=\frac{\delta_T}{r}\bigl(q^r-3q^{r/3}+2\bigr).
 \label{eq:polys}
\end{equation}
\end{enumerate}
\end{theorem}

The theorem asks for no factorization of $q^r-1$ to test a given element.
The primes and exponents in~\eqref{eq:testp} come only from $T$. A cyclic
generator is convenient in the counting proof, but is not part of the
definition or required as a distinguished label.

\section{Proof of the arithmetic properties}
\label{sec:proof}

The background facts are elementary: finite-field multiplicative groups
are cyclic; subfields containing $\F_q$ correspond to divisors of the
extension degree; and the conjugates over $\F_q$ are the Frobenius orbit.
See, for example, Conrad's account~\cite{conrad}. The new deductions from
these facts are proved in detail here.

\subsection{The coprime decomposition}

\begin{lemma}
\label{lem:coprime}
The integers in~\eqref{eq:notation} satisfy $M=HC$ and
$\gcd(H,C)=1$. In particular $\gcd(C,3)=1$ and $\gcd(C,T)=1$.
\end{lemma}
\begin{proof}
Since $B^3=q^r$, factor $B^3-1=(B-1)(B^2+B+1)$ to obtain $M=HC$.
Write $B=1+3t$. Direct calculation gives
\[
 C=1+3t+3t^2=1+(B-1)(1+t).
\]
Thus $C\equiv1\pmod{B-1}$ and $C\equiv1\pmod3$, proving
$\gcd(C,3(B-1))=1$. Finally $T\mid q-1\mid B-1$, so $C$ is also
coprime to $T$.
\end{proof}

\subsection{Why every prime-power image still has full degree}

Every proper subfield of $\F_{q^r}$ containing $\F_q$ lies in
$\F_B=\F_{q^{r/3}}$, because every proper divisor of $3^a$ divides
$3^{a-1}$. It is therefore enough to show that
$\alpha^{\ell^u}\notin\F_B$.

First suppose $\ell\ne3$. The element
\[
 z=\alpha^{M/3}=y^{m/3}
\]
is a nontrivial cube root of unity by~\eqref{eq:test3}, so
$z^{\ell^u}\ne1$. But $B-1$ divides $M/3=(B-1)C$, and hence every
$w\in\F_B^{\times}$ satisfies $w^{M/3}=1$. If
$\alpha^{\ell^u}\in\F_B$, its $M/3$-th power would be both $1$ and
$z^{\ell^u}\ne1$, a contradiction.

Now suppose $\ell=3$. Put $v=\alpha^H$. By~\eqref{eq:outside}, $v\ne1$;
by $HC=M$, $v^C=1$. Lemma~\ref{lem:coprime} shows that raising to $3^u$
is injective on the group of $C$-th roots of unity. If
$\alpha^{3^u}\in\F_B^{\times}$, then
\[
 v^{3^u}=(\alpha^{3^u})^H=1,
\]
because $B-1\mid H$. This would force $v=1$, again a contradiction.

This proves Theorem~\ref{thm:replacement}(i), including $u=0$. It uses no
lifting-the-exponent argument and no primitive-divisor theorem. The two
tests play complementary roles: the nontrivial cube root protects the
$\ell\ne3$ cases, while the nontrivial $C$-component protects the
$\ell=3$ case.

\subsection{Norm divisibility and the common power}

Let $d=\ord(y)$, which divides $m$. For a prime $p\mid m$, write
$h_p=v_p(m)$. If $p\ne3$ divides $T$, then
\[
 y^{m/p^{n_p}}\ne1
 \quad\Longleftrightarrow\quad
 d\nmid m/p^{n_p}
 \quad\Longleftrightarrow\quad
 v_p(d)>h_p-n_p.
\]
The last equivalence holds because $d\mid m$ and the displayed division
changes only the $p$-part of $m$. Since $T^2\mid m$, one has
$h_p\ge2n_p$, and therefore
\begin{equation}
 v_p(d)\ge h_p-n_p+1\ge n_p+1.
 \label{eq:margin}
\end{equation}
For $p=3$,~\eqref{eq:test3} gives $v_3(d)=h_3\ge2n_3\ge n_3$.
Thus $T\mid d$. This proves the divisibility assertion, with more than
the required margin at every prime dividing $T$.

The roots of the minimal polynomial are $\alpha^{q^i}$, $0\le i<r$.
Because $y\in\F_q$, each has $E$-th power $y^{q^i}=y$. Their product is
also $\alpha^{1+q+\cdots+q^{r-1}}=y$. If $y^x=1$, then
$d\mid x$, so $T\mid x$. This is precisely enough for the divisibility
field in \src{IsCommonRootPower}; the stronger old assertion
$d=q-1$ is unnecessary.

\subsection{Frobenius and scalar invariance}

Frobenius fixes $y$ and preserves the inequality $\alpha^H\ne1$, being
an automorphism. Hence it preserves every test.

Suppose $\lambda\in\F_q^{\times}$ and $\lambda^T=1$. Its effect on the
norm is
\[
 \Nm(\lambda\alpha)=\lambda^r y.
\]
Both $m/3$ and every $m/p^{n_p}$ occurring in the tests are multiples of
$T$. Indeed $m$ is a multiple of $T^2$, while $3\mid T$ and
$p^{n_p}\mid T$. Therefore
\[
 (\lambda^r y)^{m/3}=y^{m/3},\qquad
 (\lambda^r y)^{m/p^{n_p}}=y^{m/p^{n_p}}.
\]
Furthermore $q-1\mid B-1$, so $\lambda^H=1$ for every base-field unit
$\lambda$, and $(\lambda\alpha)^H=\alpha^H$. Normalization preserves
the actual values of the tests. No estimate for the order of a product
is needed.

\subsection{Why old primitive labels are included}

If $\alpha$ has order $M$, its norm $\alpha^E$ has order $m$. It passes
all the norm tests, because the tested exponents are proper divisors of
$m$. Also $C>1$, so $H<M$ and $\alpha^H\ne1$. Every old primitive
polynomial is therefore among the new qualifying polynomials.

The inclusion need not be strict for every parameter set. For example,
the counts agree for $(q,r,T)=(73,3,6)$. The recommendation is to weaken
the definition and simplify its proof obligations, not to assert a strict
density improvement in every individual finite field.

\section{Exact counting and the uniform bound needed by the source}
\label{sec:count}

\subsection{Counting roots by exponents}

Choose a generator $g$ of the cyclic group of order $M$ and write
$\alpha=g^e$, with $e$ taken modulo $M$. Then $g^E$ generates
$\F_q^{\times}$. The tests become
\begin{equation}
 3\nmid e,\qquad p^{n_p}\nmid e\ \ (p\mid T,\ p\ne3),
 \qquad C\nmid e.
 \label{eq:exponents}
\end{equation}
For the last equivalence, $(g^e)^H=1$ if and only if $M\mid eH$, or
$C\mid e$.

Let
\[
 D=3\prod_{\substack{p\mid T\\p\ne3}}p^{n_p}.
\]
Then $D\mid T\mid H$, and the first two tests have period $D$.
The Chinese remainder theorem shows that the fraction of residues modulo
$D$ satisfying them is exactly $\delta_T$. Thus $\delta_T H$ residues
modulo $H$ satisfy the norm tests. By Lemma~\ref{lem:coprime}, residues
modulo $M=HC$ correspond independently to residues modulo $H$ and $C$.
The last test in~\eqref{eq:exponents} permits exactly $C-1$ residues
modulo $C$. Consequently
\[
 A(q,r,T)=\delta_T H(C-1)=\delta_T(M-H)
         =\delta_T(B^3-3B+2).
\]
This proves~\eqref{eq:roots}. It is an exact count, not an asymptotic
estimate. In particular it has no dependence on the factorization of
$C$, which may be complicated.

All counted roots have degree $r$, and the tests are Frobenius invariant.
Every orbit therefore has exactly $r$ elements and corresponds to one
monic irreducible polynomial. Dividing by $r$ proves~\eqref{eq:polys}.
The factorization
\[
 B^3-3B+2=(B-1)^2(B+2)
\]
also makes positivity transparent.

\subsection{Using the source's already available exponent margin}

The general theorem does not assert that $\delta_T$ is bounded below by
$1/2$. For example, $T=6$ gives $\delta_T=1/3$. The stronger uniform
estimate uses the source's specific definition of \texttt{P.q\textsubscript{2}\ D.towerExp}:
\begin{equation}
 T=(q_0d)^{s+6}.
 \label{eq:sourceT}
\end{equation}
In the source, $3\mid q_0$, the rank is a positive power of $3$, and
$T^2\mid q-1$ at each chosen level prime. Hence the general theorem
applies, and every prime dividing $T$ has $n_p\ge6$~\cite{parameters}.

\begin{lemma}[An elementary uniform density]
\label{lem:density}
If every $n_p\ge6$, then $\delta_T\ge5/8$.
\end{lemma}
\begin{proof}
For finitely many $a_i\in[0,1]$, induction gives
$\prod_i(1-a_i)\ge1-\sum_i a_i$. For every integer $n\ge2$,
$n^6\ge16n(n-1)$, and consequently, for every finite $K$,
\[
 \sum_{n=2}^{K}\frac1{n^6}
 \le\frac1{16}\sum_{n=2}^{K}\left(\frac1{n-1}-\frac1n\right)
 \le\frac1{16}.
\]
Applying this to a finite interval containing the prime divisors of $T$
gives
\[
 \delta_T\ge\frac23\left(1-\sum_{\substack{p\mid T\\p\ne3}}p^{-6}\right)
 \ge\frac23\left(1-\frac1{16}\right)=\frac58.
\]
No infinite Euler product or prime-distribution estimate is involved.
\end{proof}

\begin{corollary}[Uniform polynomial count]
\label{cor:count}
Under~\eqref{eq:hyp} and $n_p\ge6$ for all $p\mid T$,
\begin{equation}
 \frac{P(q,r,T)}{q^r}\ge\frac5{16r}>\frac1{4r}.
 \label{eq:uniform}
\end{equation}
\end{corollary}
\begin{proof}
For $B\ge2$,
\[
 B^3-3B+2-\frac{B^3}{2}
 =\frac{(B-2)(B^2+2B-2)}2\ge0.
\]
Combine this with $q^r=B^3$,~\eqref{eq:polys}, and
Lemma~\ref{lem:density}.
\end{proof}

The constant $1/4$ deliberately leaves slack. Optimizing it would add
bookkeeping without addressing another substantial part of the proof.
The useful change is removing the dependence of this density estimate
on $r\log q$.

\subsection{A shorter counting route for the Lean implementation}
\label{sec:union}

The production proof needs only $P\ge q^r/(4r)$, so it can avoid the
Chinese remainder theorem, the exact product $\delta_T$, and exact
Frobenius-orbit counting. Rewrite the tests as
\[
 \alpha^{M/3}\ne1,\qquad
 \alpha^{M/p^{n_p}}\ne1\ (p\mid T,\ p\ne3),\qquad \alpha^H\ne1.
\]
In a cyclic group of order $M$, their respective failure sets have
cardinalities $M/3$, $M/p^{n_p}$, and $H$. The finite union bound therefore
gives, when all $n_p\ge6$,
\[
 A(q,r,T)\ge M\left(1-\frac13-
       \sum_{\substack{p\mid T\\p\ne3}}p^{-n_p}-\frac HM\right).
\]
The telescoping estimate in Lemma~\ref{lem:density} bounds the sum by
$1/16$. Also $q\ge T^2+1\ge10$, so $B\ge3$ and
\[
 \frac HM=\frac3{B^2+B+1}\le\frac14.
\]
It follows that
\begin{equation}
 A(q,r,T)\ge\frac{17}{48}M
 \ge\frac{17}{64}q^r>\frac14q^r,
 \label{eq:union}
\end{equation}
where the second inequality uses $q^r\ge4$ and hence
$M=q^r-1\ge3q^r/4$. Every counted root has a qualifying minimal
polynomial, and at most $r$ roots map to any one polynomial. Thus
$rP\ge A$ gives the same convenient bound $P\ge q^r/(4r)$.

This is the recommended counting proof for a first Lean refactor.
The current primitive-polynomial count already uses the finite-set
image/fiber lemma \src{card_le_mul_card_image_of_maps_to} with
\src{minpoly} and fibers of size at most $r$~\cite{badtypes}.
Reusing that pattern reduces the new counting obligations to elementary
power-kernel cardinalities and a finite union bound. The exact
formula~\eqref{eq:polys} can be formalized separately or retained only
as the explanatory theorem in this report.

\section{The resulting bad-type estimate}
\label{sec:weight}

Throughout this section retain~\eqref{eq:hyp} and the source-specific
assumption $n_p\ge6$ for every prime $p\mid T$.

\subsection{The comparison weight is not an exact representation mass}

Let $1<r<q$, and assign each monic polynomial of degree $r$ the uniform
comparison weight
\[
 w=\frac{1-r/q}{q^r}\ge0.
\]
The source proves that an allowed label's idempotent coefficient at the
identity is at least $w$, using Green's character degree bound. That
coefficient is generally not equal to $w$. The counting argument must
retain this distinction~\cite{badtypes}.

There are $q^r$ monic polynomials of degree $r$. By
Corollary~\ref{cor:count}, the uniform weighted count of the
nonqualifying monic polynomials is
\begin{align}
 (q^r-P)w
 &=\left(1-\frac{P}{q^r}\right)\left(1-\frac rq\right)\notag\\
 &\le1-\frac{P}{q^r}
 \le\exp\left(-\frac{P}{q^r}\right)
 \le\exp\left(-\frac1{4r}\right).
 \label{eq:weighted}
\end{align}
Equation~\eqref{eq:weighted} is \emph{not} an upper bound on the actual
Plancherel mass of all nonqualifying labels. Nor does the argument claim
that the actual mass of qualifying labels is at least $1/(4r)$.
It is a uniform weighted-count estimate used in the following cover
argument, exactly as in the existing proof.

\subsection{From counting to a tuple projector}

Let there be $f$ residue fields, all of size $q$. In each factor retain
the source's allowed-polynomial condition. Its weighted count is at least
$1-\varepsilon$, where
\begin{equation}
 \frac{1+T+2r}{q}\le\varepsilon\le1.
 \label{eq:epsilon}
\end{equation}
Define a good tuple to have all entries allowed and at least one
qualifying entry. Let $e_{\mathrm{good}}$ be the sum of the corresponding
orthogonal central tuple idempotents.

Write $\mathcal A$ for the set of all-allowed tuples and $\mathcal B$
for the set of tuples with no qualifying entry. Then
$\mathcal A\setminus\mathcal B$ is the good set. Multiplying the
elementary cardinality cover by $w^f$, applying the lower bound on every
good tuple's idempotent coefficient, and using~\eqref{eq:weighted} gives
\[
 (1-\varepsilon)^f
 \le \Re\bigl(e_{\mathrm{good}}(1)\bigr)
       +\exp\left(-\frac f{4r}\right).
\]
Here $e(1)$ means the coefficient of the identity in the group algebra.
Bernoulli's inequality $(1-\varepsilon)^f\ge1-f\varepsilon$ now yields
\begin{equation}
 \boxed{\ \Re\bigl((1-e_{\mathrm{good}})(1)\bigr)
 \le f\varepsilon+\exp\left(-\frac f{4r}\right).\ }
 \label{eq:projector}
\end{equation}
This reproduces the architecture of
\src{re_coeff_one_sub_sum_tupleIdem_le}; only the witness predicate and
the single-factor complement bound change.

\subsection{The numerical obligation that disappears}

To make the exponential term at most $e^{-L}$, the current sufficient
condition and the proposed sufficient condition are respectively
\begin{equation}
 \underbrace{2r(2r\log q+1)L\le f}_{\text{current primitive-label estimate}}
 \qquad\longrightarrow\qquad
 \underbrace{4rL\le f}_{\text{new norm-test estimate}}.
 \label{eq:f}
\end{equation}
The source applies this with $L=9\log b$. Thus the changed local
obligation is
\[
 36r\log b\le f,
\]
in place of $18r(2r\log q+1)\log b\le f$. The other error term,
$f(1+T+2r)/q$, stays the same. The caller can keep its present conclusion
$b^{-8}/2$ and the downstream near-isometry bound~\cite{spanning}.

This is a reduction of one sufficient lower bound on $f$. It does not
show that the final number of places can be reduced by that ratio:
other parameter inequalities may dominate. In particular, no improvement
of a global tower exponent, final logarithm power, build time, or total
Lean line count is asserted here.

\section{How to integrate the replacement into Lean}
\label{sec:integration}

\subsection{Preserve the output contract}

The difficult congruence consumer is
\src{IsLiftSetting.exists_isCommonRootPower_of_isPrimitivePolynomial}
in \src{RowsAndColumnsFromLabels.lean}. After normalizing the roots, it
returns a structure with four fields. The replacement supplies each one
without changing the structure~\cite{rows}.

\begin{center}
\small
\begin{tabularx}{\textwidth}{@{}p{.25\textwidth}Y@{}}
\toprule
Source field & Replacement justification \\
\midrule
\src{pow_eq} & Every root has the common norm power $y$, with exponent $E$. \\
\src{pow_q_pow_card_eq} & Every root lies in $\F_{q^r}$, so its $q^r$-th power equals itself. \\
\src{injOn_pow} & Theorem~\ref{thm:replacement}(i) gives degree $r$ for every $\ell^u$-th power. \\
\src{dvd_of_pow_eq_one} & $T\mid\ord(y)$ implies $T\mid x$ whenever $y^x=1$, hence $T\mid jx$. \\
\bottomrule
\end{tabularx}
\end{center}

The scalar in the normalization is a base-field $T$-th root of unity,
as proved in the current consumer. Section~\ref{sec:proof} proves the
new predicate invariant under exactly this operation. The separate
rank-one alternative, whose polynomial is $X-1$, must remain as it is:
the replacement theorem assumes $r\ge3$ and is not intended for that
branch.

\subsection{Definition and lemma boundaries}

Add a fresh predicate, for example
\texttt{IsQualifiedPolynomial T r c}, whose conjunction retains
\texttt{IsLabel r c} and \texttt{Irreducible c}, followed by existence of a
root satisfying the displayed tests. Keeping those conjuncts makes the
existing consumers easier to adapt even though the root tests imply
irreducibility and the degree for the counted construction.

For an implementation in an algebraic closure, one can express the norm
as the root raised to $E=1+q+\cdots+q^{r-1}$ and include the finite-field
membership statement $\alpha^{q^r}=\alpha$ in the root predicate.
Alternatively, define the predicate in a chosen finite extension and
transfer it along the standard finite-field embedding. Either route
avoids adding a chosen multiplicative generator to the public definition.

A useful order of proof obligations is:
\begin{enumerate}
\item Prove the elementary coprimality identity for $H$ and $C$.
\item Prove exclusion from the unique maximal proper subfield after
every prime-power map, and derive Frobenius injectivity.
\item Prove Frobenius invariance, scalar invariance, and norm-order
divisibility. Generalize the existing product-of-roots/common-power
identity to irreducible labels.
\item Use the finite union bound in Section~\ref{sec:union} and pass to
minimal polynomials with fibers of size at most $r$. The exact-count
theorem is useful documentation but need not be on the critical path.
\item Replace the complement estimate in \src{BadTypes.lean} by
\eqref{eq:weighted}, using the bound from the preceding step.
\item Change the distinguished witness in \src{goodTuples}; thread the
parameter $T$ through the lift-setting witness and retarget the common
root-power constructor.
\end{enumerate}

Do not silently redefine \src{IsPrimitivePolynomial}: that name has a
standard mathematical meaning. Retain it and prove inclusion in the new
class, or remove unused old lemmas later after dependency analysis.

\subsection{The audited edit sites}

The audit of explicit primitive-label references identified seven source
files. One additional file has the direct numerical call site. All paths
below are relative to \src{Lean/grh/src/}; full paths, line anchors, and
hashes appear in the accompanying machine-readable source manifest.

\begin{description}[style=nextline,font=\normalfont\bfseries]
\item[\src{New/UpperBound/Objects/LabelsAndIndexTypes.lean}]
Add the new predicate and scaling theorem; generalize the common norm
power identity; change the distinguished witness in \src{goodTuples}
and its membership theorem. Keep the allowed-label definition.
\item[\src{Lib/LocalGL/Green/LabelCharacters.lean}]
Supply the prime-power Frobenius-injectivity lemma under the weaker
conditions. The current primitive-root injectivity lemma can remain for
other uses. Green's character theory does not change.
\item[\src{New/Spanning/MainTerm/BadTypes.lean}]
Replace the primitive count and its logarithmic-totient estimate by the
new count and~\eqref{eq:weighted}. Reuse the allowed-label estimate and
the tuple cover argument.
\item[\src{New/Spanning/Multiplicities/EqualMultiplicities.lean}]
Update the distinguished-witness extraction. The trace calculation uses
only its label and irreducibility components.
\item[\src{New/Congruences/GaloisRepresentations/SelmerHypotheses.lean}]
Replace the primitive alternative in the lift-setting witness by the
new predicate; preserve the $X-1$ rank-one alternative.
\item[\src{New/Congruences/GaloisRepresentations/AttachedRepresentations.lean}]
Supply that new witness from membership in the good tuple family.
\item[\src{New/Congruences/GaloisRepresentations/RowsAndColumnsFromLabels.lean}]
Use the three new arithmetic lemmas to return the existing
\src{IsCommonRootPower} structure. Update the case split and witness
extraction that call this constructor.
\item[\src{New/Spanning/SecondMoment/SpanningTheorem.lean}]
Replace its sufficient lower bound on $f$ by $36r\log b\le f$, supplying
the source's rank and $T$ hypotheses to the new counting theorem.
\end{description}

The further uses of \src{goodTupleFamilies} in
\src{NormOfChain.lean} and \src{RowsFromEigensystems.lean} use generic
orthogonal-projector, range, or witness interfaces. The audit found no
additional demand for a root of full order $q^r-1$. The detailed audit is
included as \src{audit/lean_migration.md}.

\subsection{What must be rebuilt, and what is actually removed}

The new good projector is generally larger than the old one. The improved
count must be attached to this changed projector; it cannot simply be
inserted into an estimate about the unchanged primitive-only projector.
All uses of the definition must be rebuilt coherently. Eight substantive
edit sites are not a bound on the compiler's dependency graph; unfolding
the changed definition may require further routine elaboration edits.

The active route would no longer need
\src{totient_le_mul_card_isPrimitivePolynomial}, the
\src{le_totient_mul_log} step in this counting argument, or the demand
that the distinguished root generate all of $\F_{q^r}^{\times}$. The
direct analytic-primes import of \src{BadTypes.lean} becomes a candidate
for removal after elementary dependencies are made explicit. The
\src{Lib.Analytic.Primes} file itself still serves other prime-existence
arguments and cannot simply be deleted.

This change removes the demand for full multiplicative order in the
distinguished type label. Primitive roots used elsewhere, including the
exact-order local character and cyclotomic constructions, remain part
of the proof.

After implementing the refactor, the relevant completion gate is a
successful build of the package and inspection of the public theorem
axioms, followed by the project's statement-comparison and independent
checker workflow. This report supplies the mathematical replacement and
the audited migration plan, not results of those unperformed steps.

\section{Why this does not repeat the existing proposals}
\label{sec:novelty}

The complete proposal collection at the pinned revision has five
directories, containing 47 files: 42 text files and five PDFs. Every main
LaTeX source, included section, README, and formalization plan was read.
The PDF arguments were reviewed through their complete corresponding
LaTeX sources. All 47 Git blobs were checked against the pinned checkout.
The detailed mechanism-by-mechanism audit and file inventory accompany
this report~\cite{proposals}.

\begin{center}
\small
\begin{tabularx}{\textwidth}{@{}p{.31\textwidth}Y@{}}
\toprule
Existing directory & Improvements treated as already available \\
\midrule
\src{centered_theta_lift_simplification} & Centered covariance, direct rank increment, adaptive completion, smaller translate budgets, optional random contractions. \\
\src{proof_simplifications} & Large-sieve or Fourier good-twist selection, smoothed Perron inversion at the distant cutoff, direct ramified spanning. \\
\src{rh-simplifications} & Sharp fourth-order rank deficit, torsion interpolation, mollifier/large-sieve twist selection, endpoint assembly. \\
\src{rh_sparse_theta_proof} & Reciprocal-cutoff convexity, sparse grids, averaged cutoffs, Gaussian moments, weaker prime-subfamily requirements. \\
\src{shorter-theta-lift-proof} & Correlations and Poisson completion, mollification, disk zero detection, centered spanning, PNT cutoffs and interpolation. \\
\bottomrule
\end{tabularx}
\end{center}

The audit also excluded the future directions explicitly proposed in
those documents: a unified rank--height budget, index charged by actual
rank, partial/exterior determinants, weaker orbital estimates at a
chosen tolerance, shared fixed-chain data, favorable prime subfamilies,
and direct averaged volumes. Rephrasing one of those directions would
not satisfy the user's request.

None of the audited proposals replaces the primitive witness in the
finite good-type projector or proves the norm-order and prime-power
Frobenius properties under a weaker condition. None supplies the count
\eqref{eq:polys}. That is the distinct mathematical contribution here.
It remains useful after the proposed spanning and analytic changes are
assumed implemented: the replacement preserves the local output contract
and improves an earlier projector estimate on which those arguments can
continue to rely.

This is a nonduplication claim relative to the requested repository
collection at the recorded commit. It is not a claim of global literature
priority for the elementary finite-field ideas.

\section{Examples, counterexamples, and executable checks}
\label{sec:checks}

\subsection{A qualifying polynomial that is not primitive}

Over $\F_{19}$ take
\[
 c(X)=X^3+X+14,\qquad r=3,\qquad T=3.
\]
It has no root in $\F_{19}$, hence is irreducible. For
$\alpha=X\bmod c$, exact arithmetic gives
\[
 \ord(\alpha)=3429,\qquad 19^3-1=6858,
 \qquad\Nm(\alpha)=5,\qquad\ord(5)=9.
\]
Also
\[
 \alpha^{54}=16+9\alpha+12\alpha^2\ne1.
\]
The norm test is $5^6\ne1$ and the exclusion exponent is
$3(19-1)=54$. Thus $c$ qualifies although its root has only half the
full multiplicative order. This example concretely separates the new
condition from primitivity.

For these parameters the exact count is
\[
 P(19,3,3)=\frac{2}{9}(19^3-3\cdot19+2)=1512,
\]
whereas there are $\varphi(6858)/3=756$ primitive cubics. An independent
enumeration of all $19^3=6859$ monic cubics found 2280 irreducible cubics,
1512 qualifying cubics, and 756 primitive cubics.

\begin{center}
\small
\begin{tabular}{@{}rrrrrr@{}}
\toprule
$q$ & $r$ & $T$ & $\delta_T$ & Qualifying polynomials & Primitive polynomials \\
\midrule
19 & 3 & 3 & $2/3$ & 1,512 & 756 \\
37 & 3 & 6 & $1/3$ & 5,616 & 4,752 \\
73 & 3 & 6 & $1/3$ & 43,200 & 43,200 \\
109 & 3 & 6 & $1/3$ & 143,856 & 123,120 \\
64 & 3 & 3 & $2/3$ & 58,212 & 46,656 \\
19 & 9 & 3 & $2/3$ & 23,902,790,904 & 11,834,224,416 \\
\bottomrule
\end{tabular}
\end{center}

These small moduli test the general theorem. They do not satisfy the
source-specific condition $v_p(T)\ge6$ and are not numerical tests of
that corollary's parameter regime.

\subsection{Why the hypotheses cannot simply be dropped}

Let $g$ generate $\F_{19^3}^{\times}$, of order $6858=54\cdot127$.
\begin{enumerate}
\item \textbf{Omitting the subgroup exclusion fails.} The element
$g^{127}$ has order $54$ and primitive norm of order $18$, but
$(g^{127})^3$ has order $18$ and lies in $\F_{19}$. The degree condition
fails for $\ell=3$.
\item \textbf{Omitting the full $3$-part of the norm fails.} The element
$g^{27}$ has order $254$ and passes the subgroup exclusion, but its
$127$-th power is $-1\in\F_{19}$. A condition only about being outside
the exceptional subgroup does not suffice.
\item \textbf{Merely imposing $T\mid\ord(y)$ is not stable under the
required scalars.} For $(q,r,T)=(37,3,6)$ and a generator $g$ of
$\F_{37^3}^{\times}$, the norm of $g^2$ has order $18$. Multiplying by
$-1$, which is a $T$-th root of unity, changes the norm order to $9$,
which is not divisible by $6$. The margin encoded in the power tests
avoids this cancellation.
\item \textbf{The rank hypothesis is substantive.} If $r=6$, $q=19$,
$T=3$, and $g$ generates $\F_{19^6}^{\times}$, then $\alpha=g^{1715}$
has order $27432$ and degree $6$, has primitive norm, and satisfies
$\alpha^{3(19^2-1)}\ne1$. Nevertheless $\alpha^8$ has order $3429$
and degree $3$. For a rank with another prime divisor, the field of
degree $r/3$ no longer contains every proper subfield.
\end{enumerate}

These examples explain the design of the predicate. It should not be
weakened to ``an irreducible label with a sufficiently large norm'' without
rechecking both normalization and all prime-power maps.

\subsection{What the included programs verify}

\src{code/check_order_tests.py} uses only the Python standard library
and exact integer or rational arithmetic. It checks the count by
exhausting element exponents in the five rank-three cases in the table;
tests Frobenius, scalar invariance, norm divisibility and prime-power
degrees; and independently implements arithmetic in cubic quotient
fields to enumerate all cubics over $\F_{19}$. Larger structural tests
use explicitly recorded deterministic samples. The rank-nine root count
is not checked by enumerating its hundreds of billions of elements.

\src{code/check_by_orders.py} is an independent checker using a different
counting route: it sums $\varphi(d)$ over qualifying element orders
$d\mid q^r-1$. It also exhausts selected small cyclic groups and their
allowed scalars. Its ``power'' mode is the predicate recommended here;
the other two modes compare neighboring norm predicates and are
auxiliary checks, not additional recommendations.

The source checker verifies the pinned commit, source-file hashes,
declaration anchors, and all proposal blobs. The saved JSON outputs and
guide-check output are supplied alongside the scripts. Finite checks
guard against algebraic mistakes and confirm the displayed examples;
the general proof is the argument in
Sections~\ref{sec:theorem}--\ref{sec:weight}.

\section{Assessment and next implementation step}

The strongest reason to adopt this change is the match between a smaller
hypothesis and the actual contract. The source needs one irreducible
label with stable prime-power Frobenius orbits and enough common-power
order. Requiring a generator of a much larger cyclic group hides this
structure and brings an avoidable totient estimate into the density
argument. Norm power tests make the required behavior explicit, and the
coprime factorization $M=HC$ makes the count elementary.

The proposed change retains the rank $3^a$ and the square divisibility
$T^2\mid q-1$. It retains the allowed labels, Green's character estimates,
the rank-one branch, the local Galois and Selmer machinery after
\src{IsCommonRootPower}, and the theta-lift volume comparison. Those are
substantial limitations on the size of the simplification. They also
make the replacement sufficiently local to be reviewable and compatible
with the existing proposals.

A sensible first implementation is to keep every global parameter at
its present value, prove the new local lemmas, change the good projector,
and recover the existing public bounds. Only after that build succeeds
should one revisit the global parameter ledger to determine whether
the weaker bound on $f$ removes any other dominant constraint. That
second optimization is not assumed in this report's conclusions.

\begin{keybox}
\textbf{Concrete result.} Theorem~\ref{thm:replacement} proves the new
finite-field replacement, Corollary~\ref{cor:count} supplies the uniform
density for the repository's parameters, and~\eqref{eq:projector}
supplies the improved projector estimate. The source audit identifies
the contracts that must change and finds no downstream use requiring
the discarded full multiplicative order. A complete Lean implementation
and rebuild remain to be done.
\end{keybox}

\appendix
\section{Reproducing the report and checks}

From the extracted archive, the main exact checks can be run as follows:
\begin{verbatim}
python3 code/check_order_tests.py --output checks/exact_checks.json
python3 code/check_by_orders.py --output checks/order_checks.json
\end{verbatim}
To verify the source audit, obtain a checkout of ProveIt at the recorded
commit and run:
\begin{verbatim}
python3 code/check_source_snapshot.py /path/to/ProveIt
\end{verbatim}
The script is read-only with respect to that checkout. It verifies the
source and proposal manifests included here; it does not download
dependencies or build Lean.

The self-contained LaTeX source uses standard TeX Live packages. Run the
following command three times to settle the contents and cross-references:
\begin{verbatim}
pdflatex -halt-on-error order_tests.tex
\end{verbatim}
The supplied PDF was compiled and visually checked. The archive also
contains file checksums, source metadata, and the recorded verification status.

\section{Notation and source dictionary}
\begin{center}
\small
\begin{tabularx}{\textwidth}{@{}p{.19\textwidth}Y@{}}
\toprule
This report & Meaning and source counterpart \\
\midrule
$q$ & Size of the residue field at a chosen place; in the final caller, \texttt{D.levelPrime i}. \\
$r$ & Rank; \src{D.r} or \texttt{card D.RowIndex}; a positive power of $3$. \\
$T$ & Local-character order \texttt{P.q\textsubscript{2}\ D.towerExp}; called $M$ in generic good-tuple definitions and $N$ in some source docstrings and check programs. \\
$s$ & Tower exponent in $T=(q_0d)^{s+6}$. \\
$f$ & Number of residue-field factors in the tuple; \src{D.f}. \\
$b$ & Weight parameter, with \src{D.logB} equal to $\log b$. \\
$B$ & Only in this report's finite-field calculation, $q^{r/3}$; unrelated to the source's uppercase log notation. \\
$H,C$ & Coprime factors $3(B-1)$ and $(B^2+B+1)/3$ of $q^r-1$. \\
$y$ & The common norm of the roots, also their product. \\
$P(q,r,T)$ & Number of qualifying monic polynomials; not the source's parameter structure \src{P}. \\
\bottomrule
\end{tabularx}
\end{center}

\begin{thebibliography}{20}
\small
\bibitem{guide} ProveIt, \emph{How theta lifts locate the zeros of Dirichlet
$L$-functions}, guide and source map, pinned repository revision.
\href{\repobase/docs/guide/theta_lifts_and_dirichlet_zeros.tex}{Guide source}.
\bibitem{mainineq} ProveIt,
\sourcefile{New/Reduction/MainInequality.lean}{\emph{MainInequality.lean}},
especially \src{star_inequality}.
\bibitem{onesided} ProveIt,
\sourcefile{New/Reduction/OneSidedBound.lean}{\emph{OneSidedBound.lean}},
especially \src{eventually_primeRemainder_le_rpow} and \src{grh}.
\bibitem{labels} ProveIt,
\sourcefile{New/UpperBound/Objects/LabelsAndIndexTypes.lean}{\emph{LabelsAndIndexTypes.lean}},
definitions and lemmas for primitive polynomials, allowed polynomials,
and good tuples.
\bibitem{multiplicities} ProveIt,
\sourcefile{New/Spanning/Multiplicities/EqualMultiplicities.lean}{\emph{EqualMultiplicities.lean}},
distinguished-witness extraction and \src{trace_mul_tupleIdem_eq}.
\bibitem{rows} ProveIt,
\sourcefile{New/Congruences/GaloisRepresentations/RowsAndColumnsFromLabels.lean}{\emph{RowsAndColumnsFromLabels.lean}},
\src{IsCommonRootPower} and its constructors.
\bibitem{badtypes} ProveIt,
\sourcefile{New/Spanning/MainTerm/BadTypes.lean}{\emph{BadTypes.lean}},
the primitive count, weighted complement bound, and tuple cover argument.
\bibitem{spanning} ProveIt,
\sourcefile{New/Spanning/SecondMoment/SpanningTheorem.lean}{\emph{SpanningTheorem.lean}},
\src{re_coeff_one_sub_goodIdempotentAbove_le} and \src{bad_types_rare}.
\bibitem{parameters} ProveIt,
\sourcefile{New/Reduction/Data/FieldsAndPlaces.lean}{\emph{FieldsAndPlaces.lean}}
and
\sourcefile{New/Reduction/Data/AdmissibleData.lean}{\emph{AdmissibleData.lean}},
the definition of \texttt{q\textsubscript{2}}, the rank power, and square divisibility at
level primes.
\bibitem{proposals} ProveIt, five existing
\href{https://github.com/VladimirReshetnikov/ProveIt/tree/89d12dd40424f1217fd818d84877b3da52396a3b/Analysis/GeneralizedRiemannHypothesis/docs/proposals}{proposal directories}.
See \src{audit/existing_proposals.md} and its complete blob inventory
in this archive for individual documents and precise overlap locations.
\bibitem{sutherland} Andrew Sutherland,
\href{https://math.mit.edu/~drew/shortgrh/}{\emph{Verification files for the
shorter Lean proof of RH and of GRH for Dirichlet $L$-functions}},
page updated 27 November 2026; consulted 27 November 2026.
\bibitem{conrad} Keith Conrad,
\href{https://kconrad.math.uconn.edu/blurbs/galoistheory/finitefields.pdf}{\emph{Finite Fields}},
background on cyclic multiplicative groups, subfields, Frobenius, and
finite-field norms.
\end{thebibliography}
\end{document}
